\section*{الفصل \ref{ch:g4:clouds-rain-snow} --- السحب والمطر والثلج}

\begin{solution}{exo:g4:clouds-rain-snow:1}
تتكون \omterm{def:g4:clouds-rain-snow:cloud}{السحابة} من عدد كبير جدًا من قطيرات الماء الصغيرة جدًا، أو من قطع
الجليد الصغيرة جدًا، المعلقة في الهواء. \omterm{def:g4:clouds-rain-snow:cloud}{والضباب} \omterm{def:g4:clouds-rain-snow:cloud}{سحابة} تلامس الأرض.
\end{solution}

\begin{solution}{exo:g4:clouds-rain-snow:2}
(أ) \omterm{def:g4:clouds-rain-snow:cloud-families}{طبقي}. (ب) \omterm{def:g4:clouds-rain-snow:cloud-families}{سمحاق}. (ج) \omterm{def:g4:clouds-rain-snow:cloud-families}{ركام}. (د) \omterm{def:g4:clouds-rain-snow:cloud-families}{ركام مزني}.
\end{solution}

\begin{solution}{exo:g4:clouds-rain-snow:3}
(أ) \omterm{def:g4:clouds-rain-snow:precipitation}{برد}. (ب) \omterm{def:g4:clouds-rain-snow:precipitation}{مطر}. (ج) \omterm{def:g4:clouds-rain-snow:precipitation}{ثلج}.
\end{solution}

\begin{solution}{exo:g4:clouds-rain-snow:4}
(أ) \omterm{def:g4:clouds-rain-snow:rain-gauge}{مقياس المطر}. (ب) مقياس الحرارة. (ج) \omterm{def:g4:clouds-rain-snow:wind-vane}{دوارة الرياح}.
\end{solution}

\begin{solution}{exo:g4:clouds-rain-snow:5}
سقط 17 مليمترًا من \omterm{def:g4:clouds-rain-snow:precipitation}{المطر}. $20 - 17 = 3$: 3 مليمترات أخرى
توصل الماء إلى العلامة 20.
\end{solution}

\begin{solution}{exo:g4:clouds-rain-snow:6}
أكثر الأيام مطرًا الأربعاء (12 مليمترًا). ولم تمطر يوم الاثنين
ولا يوم الجمعة. وكان في الخميس 7 مليمترات.
\end{solution}

\begin{solution}{exo:g4:clouds-rain-snow:7}
ندف \omterm{def:g4:clouds-rain-snow:precipitation}{الثلج} تتكون عاليًا في \omterm{def:g4:clouds-rain-snow:cloud}{السحابة} الباردة. وفي أعلى، في القرية الجبلية،
الهواء متجمد طوال الطريق إلى الأسفل، فتصل ثلجًا. وفي الأسفل، في \omterm{def:g1:landscapes-around-us:landform}{الوادي}،
الهواء أدفأ من \SI{0}{\celsius}: فتذوب الندف وهي تعبره نازلة،
وتصل إلى الأرض مطرًا.
\end{solution}

\begin{solution}{exo:g4:clouds-rain-snow:8}
لمخطط الأعمدة سبعة أعمدة ارتفاعاتها 3 و0 و1 و9 و6 و0 و1 (لا عمود يوم
الثلاثاء ويوم السبت). وفي الأسبوع سقط $3 + 0 + 1 + 9 + 6 + 0 + 1 = 20$
مليمترًا. والخميس جاء بتسعة مليمترات من العشرين: أي $\frac{9}{20}$
من \omterm{def:g4:clouds-rain-snow:precipitation}{مطر} الأسبوع.
\end{solution}

\begin{solution}{exo:g4:clouds-rain-snow:9}
الجدران والأشجار تحجب \omterm{def:g4:clouds-rain-snow:precipitation}{المطر} عن الأرض القريبة منها، ولا سيما حين تدفع الريح
\omterm{def:g4:clouds-rain-snow:precipitation}{المطر} جانبيًا، وقد تقطر القطرات من الأوراق في المقياس: فقربها يجمع
المقياس أقل مما يجب، أو أكثر. ومقياس الحرارة يجب أن يقيس درجة حرارة
الهواء؛ وفي الشمس الساطعة تسخنه أشعة الشمس نفسها فيقرأ قيمة أعلى مما يجب.
والصندوق الأبيض ذو الشرائح يحجب الشمس ويسمح بمرور
الهواء.
\end{solution}

\begin{solution}{exo:g4:clouds-rain-snow:10}
\omterm{def:g4:clouds-rain-snow:wind-vane}{دوارة الرياح} تشير إلى حيث تأتي الريح: فالريح تأتي من
الغرب. وهي تذهب نحو الشرق.
\end{solution}

\begin{solution}{exo:g4:clouds-rain-snow:11}
مليمتر \omterm{def:g4:clouds-rain-snow:precipitation}{المطر} على متر مربع واحد لتر واحد. وعلى متر مربع واحد تعطي 8
مليمترات 8 لترات. وعلى 500 متر مربع:
$500 \times 8 = 4000$ لتر. وبمرشات سعة كل منها 10 لترات:
$4000 \div 10 = 400$ مرش.
\end{solution}
